%% =====================================================================
%%  B3-T-28-02 -- what B3 contributes to "The Goal Line"
%%  -------------------------------------------------------------------
%%  LAYER A: APPEND-ONLY. Written once at the entry's write, then left
%%  alone. If the source has more to say later, add a bullet -- do not
%%  rewrite. Material found ONLY here goes in the `unique` slot with
%%  @unique set to yes, which makes it undeletable (rule R10).
%% =====================================================================
%% @kind: source
%% @id: B3-T-28-02
%% @book: B3
%% @topic: T-28-02
%% @locator: Law 47, pp. 410--418
%% @status: distilled
%% @unique: no
%% @integrated: yes
%% @updated: 2026-09-19

\dossier{B3}{T-28-02}{\bookshort{B3}}{Law 47, pp. 410--418}

\begin{distilled}
  \item The moment of victory is often the moment of greatest peril: arrogance and overconfidence push past the goal aimed at, and going too far makes more enemies than the win defeated.
  \item Do not allow success to go to the head; there is no substitute for strategy and careful planning: set a goal, and on reaching it, stop.
  \item Cyrus as the transgression case: king of the world after Babylon, warned by Queen Tomyris to be content and keep his forces whole, he crossed into her country for a final conquest and --- by her oath --- was given more blood than he could drink.
  \item The keys add the internal case: those who push past the mark are often proving dedication to a master or an audience --- the intoxication is social, not just private.
  \item The reversal: Machiavelli's fork --- either destroy a man entirely or leave him entirely alone; the law prices stopping, not surrender.
\end{distilled}
